\documentclass[a4paper,12pt]{article}
\usepackage[utf8]{inputenc}
\usepackage[ngerman]{babel}
\usepackage{amsmath}
\usepackage{geometry}
\geometry{margin=2.5cm}
\usepackage{booktabs}
\usepackage{parskip}
\usepackage{tikz}
\usepackage{href-ul}
\hypersetup{
	colorlinks=true,
	linkcolor=blue,
	urlcolor=blue}
\date{}

\begin{document}
	
		\begin{tikzpicture}(10,3)
		\draw[ultra thick](2,0) --(10,0) -- (10,1.3) --(2,1.3) -- (2,0);
		\draw[fill=black](2,0)-- (0,.6) -- (2,.6) -- (2,0);
		\node at (6,.6) {\large Lösungsblatt von \href{https://www.okuyakl.de}{www.okuyakl.de}};
	\end{tikzpicture}
	\vspace{0.5 cm}
	\section*{1. Der hungrige Hamster auf dem Mond}
	
	\begin{itemize}
		\item a) \( m = 0{,}150\,\text{kg} \), \( g_\text{Erde} = 9{,}81\,\frac{\text{m}}{\text{s}^2} \)
		
		\[
		F = 0{,}150 \cdot 9{,}81 = \boxed{1{,}47\,\text{N}}
		\]
		
		\item b) \( g_\text{Mond} = 1{,}6 \)
		
		\[
		F = 0{,}150 \cdot 1{,}6 = \boxed{0{,}24\,\text{N}}
		\]
		
		\item c) Unterschied in Prozent:
		
		\[
		\frac{1{,}47 - 0{,}24}{1{,}47} \cdot 100 \approx \boxed{83{,}7\,\%}
		\]
	\end{itemize}
	
	\hrulefill
	
	\section*{2. Die Yoga-Einlage des Elefanten}
	
	\begin{itemize}
		\item a) \( F = 3000 \cdot 9{,}81 = \boxed{29430\,\text{N}} \)
		
		\item b) Umstellen: \( m = \frac{F}{g} = \frac{33\,000}{10} = \boxed{3300\,\text{kg}} \)
	\end{itemize}
	
	\hrulefill
	
	\section*{3. Der Heliumballon in der Sauna}
	
	\begin{itemize}
		\item \( m = 0{,}050\,\text{kg} \), \( g = 9{,}81 \)
		
		\[
		F = 0{,}050 \cdot 9{,}81 = \boxed{0{,}49\,\text{N}}
		\]
	\end{itemize}
	
	\hrulefill
	
	\section*{4. Die intergalaktische Kofferwaage}
	
	\begin{itemize}
		\item \( m = 25\,\text{kg} \)
		
		\begin{itemize}
			\item Erde: \( 25 \cdot 9{,}81 = \boxed{245{,}25\,\text{N}} \)
			\item Mars: \( 25 \cdot 3{,}71 = \boxed{92{,}75\,\text{N}} \)
			\item Jupiter: \( 25 \cdot 24{,}79 = \boxed{619{,}75\,\text{N}} \)
			\item Uranus: \( 25 \cdot 8{,}69 = \boxed{217{,}25\,\text{N}} \)
		\end{itemize}
		
		\item b) Am schwersten: \boxed{Jupiter}
		
	\end{itemize}
	
	\hrulefill
	
	\section*{5. Das schwer verliebte Hufeisen}
	
	\begin{itemize}
		\item a) \( m = 2\,\text{kg} \), \( g = 9{,}81 \), also:
		
		\[
		F = 2 \cdot 9{,}81 = \boxed{19{,}62\,\text{N}}
		\]
		
		\item b) Die Magnetkraft muss mindestens \boxed{19{,}62\,\text{N}} betragen, sonst fällt das Eisen runter.
		
		\item c) Liebe ist messbar, wenn sie eine Anziehungskraft hat. In diesem Fall: \textbf{19,62 N romantische Spannung}.
	\end{itemize}
	\begin{center}
		\includegraphics[width=7 cm]{../../viecher/endcomic.pdf}
		
		Hier geht es zurück zum \href{https://www.okuyakl.de/phys/p7fmgL058/aa058.pdf}{Aufgabenblatt}
	\end{center}
\end{document}
